\chapter{Visualization of Experimental Results}
\lhead{CHAPTER 5. VISUALIZATION OF EXPERIMENTAL RESULTS}
\label{chap:visualization}

This chapter provides figure placeholders corresponding to the main trends discussed in Chapter~\ref{chap:results}. Each plot is intended to make the quantitative findings easier for non-technical stakeholders to interpret.

\begin{figure}[H]
\centering
\fbox{\parbox[c][5cm][c]{0.9\textwidth}{\centering Placeholder: Reward convergence curve}}
\caption{Reward convergence curve showing training progress over episodes for the proposed distributional RL scheduler.}
\label{fig:reward_convergence}
\end{figure}

Figure~\ref{fig:reward_convergence} is expected to show gradual stabilization of episodic reward, indicating that the C51-based policy learns more consistent scheduling behavior over time.

\begin{figure}[H]
\centering
\fbox{\parbox[c][5cm][c]{0.9\textwidth}{\centering Placeholder: Makespan comparison bar chart}}
\caption{Makespan comparison across FIFO, Fair Scheduler, Standard DQN, Verma et al. (2025), and the proposed method.}
\label{fig:makespan_comparison}
\end{figure}

Figure~\ref{fig:makespan_comparison} should clearly illustrate reduced total completion time for the integrated extension model.

\begin{figure}[H]
\centering
\fbox{\parbox[c][5cm][c]{0.9\textwidth}{\centering Placeholder: Cost distribution under stochastic pricing}}
\caption{Distribution of scheduling cost under stochastic VM pricing, comparing static and spot-aware policies.}
\label{fig:cost_distribution}
\end{figure}

Figure~\ref{fig:cost_distribution} highlights the economic benefit of stochastic price-aware scheduling, where lower median and tighter spread indicate better cost control.

\begin{figure}[H]
\centering
\fbox{\parbox[c][5cm][c]{0.9\textwidth}{\centering Placeholder: Job batching throughput plot}}
\caption{Throughput improvement with increasing batch size in job batching pre-processing.}
\label{fig:batching_throughput}
\end{figure}

Figure~\ref{fig:batching_throughput} is expected to show that throughput rises as batching reduces repeated scheduling overhead.

\begin{figure}[H]
\centering
\fbox{\parbox[c][5cm][c]{0.9\textwidth}{\centering Placeholder: Transfer learning convergence curves}}
\caption{Training curve comparison for scratch, transfer-frozen, and transfer-finetune settings under workload shift.}
\label{fig:transfer_curves}
\end{figure}

Figure~\ref{fig:transfer_curves} should demonstrate faster cost convergence under transfer settings and clarify the return-versus-cost trade-off observed in Table~\ref{tab:ext4_results}.

In summary, the five visualizations create a stakeholder-friendly narrative: learning stability, faster completion, lower cost, lower overhead, and faster adaptation.
